% Figure for Classical Mechanics §B3.1 (decomposition about the centre of mass: r_i = R + r_i').
\begin{tikzpicture}[>={Stealth[length=2.4mm]}, font=\small]
  % body outline (a blob of particles)
  \draw[black!55, thin, fill=black!4] plot[smooth cycle, tension=0.8] coordinates {(3.2,1.6) (4.6,1.3) (6.0,1.9) (6.4,3.1) (5.6,4.1) (4.1,4.0) (3.0,3.0)};
  % particles
  \foreach \p in {(3.6,2.1),(4.75,3.75),(5.4,3.6),(5.9,2.4),(5.3,2.05),(4.9,1.8)} \fill[black!60] \p circle (1.6pt);
  % origin
  \fill (0,0) circle (2pt) node[below left] {$O$};
  % centre of mass
  \coordinate (C) at (4.7,2.75);
  \fill[figR] (C) circle (2.4pt);
  \node[figR, below right, xshift=-2pt] at (C) {CM};
  % particle i
  \coordinate (P) at (3.55,3.45);
  \fill[figB] (P) circle (2.2pt);
  \node[figB, above left, yshift=1pt] at (P) {$m_i$};
  % vectors
  \draw[->, very thick, figR] (0,0) -- (C) node[pos=0.55, below right] {$\vec R$};
  \draw[->, very thick, figB] (0,0) -- (P) node[pos=0.5, above left] {$\vec r_i$};
  \draw[->, very thick, figG] (C) -- (P) node[pos=0.45, below left, xshift=1pt, yshift=1pt] {$\vec r_i^{\,\prime}$};
  % velocity of CM and of particle
  \draw[->, thick, figR, dashed] (C) -- ++(1.5,-0.6) node[right] {$\vec V$};
  \draw[->, thick, figB, dashed] (P) -- ++(0.5,1.35) node[above] {$\vec v_i = \vec V + \vec v_i^{\,\prime}$};
  % caption-like note
  \node[font=\scriptsize, black!70, align=left, anchor=west] at (3.6,0.75) {$\sum_i m_i\vec r_i^{\,\prime} = \vec 0,\ \ \sum_i m_i\vec v_i^{\,\prime} = \vec 0$};
\end{tikzpicture}
